\begin{center}
	\textbf{\heiti \large 序 言}
\end{center}


我怀着十分高兴的心情向读者介绍和推荐邹应同志编写的 《数学分析》.

时代在前进, 数学与其他自然科学技术的飞速发展对高等学校的数学教学提出了新的更高的要求. 教材内容的现代化就是数学教学适应时代发展的一个十分重要的方面. 《数学分析》的作者对传统的分析教材从体系、结构、选材及处理手法上作了大胆的有益的改革探索. 可以说这是目前国内难度最大、最接近现代数学的一本数学分析教科书。无疑, 此书的出版, 对数学教学水平的提高,对大学生的数学素质的培养将发挥重要的作用.

本书的主要特色可归纳为以下几点.

1. 起点高。

作者一改传统分析教材的模式, 首先通过介绍集合论的基本概念及其性质,着重研究了集合的“关系”,并利用“关系”的概念严格定义了映射, 利用有理数Cauchy 序列等价类定义了实数, 继而详细研究了实数的连续性, 这就不仅为一元实值函数的微积分学奠定了一个严密的数域, 而且使整个教材建立在一个高起点之上, 为以后各章的深入展开铺平了道路.

2. 注重分析、代数、几何三部分知识的相互渗透,相互联系.

现代数学发展的特点就是各个数学分支相互沟通、关联, 学科界限越来越小. 本书虽是一本数学分析基础课教材, 但作者把握了现代数学发展的这一新特点,将大学一、二年级代数中的群、同态、同构、线性映射、矩阵, 几何中的曲线曲面等概念、 知识与分析揉合在一起,并且是一种有机的结合, 这就把学生所学的知识融为一个整体,可以培养学生综合应用所学知识分析问题和解决问题的能力。

3. 与高年级后继课程紧密结合。

作者在介绍经典微积分学理论的同时,将现代数学中的许多概念、理论恰到好处地介绍到分析教材中。例如,通过介绍度量空间、点集拓扑的概念及其性质将传统的分析教材的二维、三维空间, 二元、三元函数的极限与连续性用新的更高的观点作了统一处理,又例如积分部分的处理手法也颇具特色,作者用阶梯函数序列对定义了函数的Riemann可积性及其Riemann 积分, 并用积分定义了集合的可测性及其Jordan测度；对经典的第一型、 第二型曲面积分,作者通过介绍微分形式与 \({\mathbf{R}}^{n}\) 的子流形,用流形上的积分的形式给予了全新的表述, 通过证明流形上的 Stokes公式而将经典的Newton-Leibniz公式, 平面Green公式, 空间Stokes公式与 Ostrogradski 公式作为它的特例进行了统一处理. 这不仅使教材内容本身现代化了,而且为学生在高年级学习拓扑、泛函分析、实变函数、微分几何等课程提供了一个现实模型, 为学生学习这些课程打下了一个良好的基础. 可以说, 本书起到了沟通大学两个学习阶段的桥梁作用.

4. 注重新概念、新理论的深入浅出的引入.

教材内容要现代化,要与后继课程衔接,这就涉及到新概念, 新理论的介绍. 作者不是简单地罗列新概念, 新结论, 而是通过分析现实例子,找出可推广的因素及其实质所在,自然地引出了新概念, 新理论. 例如作者从映射的角度通过分析实数的绝对值与数乘的性质而引入了抽象的度量空间与内积空间的概念. 通过分析导数的定义而引入了用线性映射定义映射微分的概念. 这为今后在流形的切空间上引进微分映射 (即切映射) 提供了实例和基础. 作者在这里不仅让学生对新概念 “知其然”,而且使学生 “知其所以然”.

5. 配备有开发学生智力的习题.

本书不仅有大量的让学生掌握基本概念与基本技能的基础性题目,而且还配备有一定数量的扩充教材内容的习题,这些题目无疑对开发学生智力,培养学生独立理解新概念,新结论,独立分析问题和解决问题的能力,提高学生的学习积极性及培养学生的创造性是十分有益的.

数学分析教材改革的难度很大,本书作者邹应同志在教材改革上迈出了可喜的一步。我深信, 本书的出版将会对数学分析教材和教学的改革以及数学分析的现代化起很大的推动作用.
\begin{flushright}
	徐森林\\
	1992年10月1日于合肥
\end{flushright}

\newpage
\begin{center}
	\textbf{\heiti \large 编 者 的 话}
\end{center}

本书是我为武汉大学数学系中法数学与计算机科学中心的数学试验班及国家理科基础科学研究和教学人才培养基地武汉大学数学基地班的一、二年级学生讲授数学分析课而编写的。

1980 年根据中法文化交流协议,在武汉大学数学系成立中法数学中心并创办中法数学试验班。试验班按照法国的教学大纲, 使用法国的教材,由中法双方教师共同进行教学. 经过几年的教学改革试验, 到1985年, 试验班进入了一、二年级的教学由中方教师承担,学制由五年改为四年的新的发展阶段。为了适应这一新形势的要求, 试验班决定开始组织编写自己的能与法国教师的高年级的数学教学相衔接的基础课教材。我的《数学分析》就是在这一形势下产生的.

书的初稿形成于1985年. 在中法试验班经过两年的试教后, 我对初稿作了比较大的修改,形成了第二稿,并先后由孙道椿、 朱赋鎏、王小林先生继续在试验班试用. 1990年在高等教育出版社文小西先生的大力支持和热情关心下,将书稿送中国科学技术大学徐森林教授审查。徐森林教授对送审稿进行了认真的审阅, 对书稿的结构、特点作了充分肯定,并提出了进一步修改和加强的宝贵意见. 我根据他所提的建议对书稿进行了全面的调整、修改、充实并重写了部分章节,形成第三稿,由我和王小林先生在数学基地班为 91 级、 92 级及 93 级学生试讲. 在出版之前根据试教过程中发现的问题及责任编辑文小西先生所提的意见又进行了最后一次修改,最终定稿形成现在与读者见面的《数学分析》一书。

从本书的结构来看大体可以分为四大块:

第一块: 第1-9章的一元函数的微积分.

第二块: 第10章的度量空间.

第三块: 第11-14章的数项与函数项级数.

第四块: 第15-21章的多元函数微积分.

在一元函数微积分部分, 考虑到现行高三数学教材已经对集合、映射的基本概念有所介绍,因此我把书的起点定得高一些, 以集合及其基本性质为基础,利用集合的关系、逆关系严格定义了映射、逆映射与复合映射等概念, 继而比较顺利地介绍了 Cantor 的实数理论. 教学实践表明, 只要集合的关系、等价关系、 序关系三个概念讲得透彻,学生对较早介绍严格的实数理论是可以接受的.

与传统的分析教材讲法不一样的是, 我把积分论放在导数的前面讲, 主要出发点 (这两部分内容讲的先后次序影响不大) 是通过以积分函数定义基本初等函数向学生介绍基本初等函数的不同定义方法. 对于一元实值函数的Riemann积分 (多元函数的重积分也是如此), 我没有采用通常的达布上和、下和的定义方法, 而是用阶梯函数逼近的方法, 用阶梯函数序列对定义函数的 Riemann可积性及其 Riemann 积分. 这样处理有一个好处: 当把绝对值换成一般范数时,可以把这种积分理论推广到定义在完备赋范向量空间即Banach空间上的映射上去,这就为学生以后学习其它类型的积分理论打下一个基础.

对于函数的导数部分, 除了介绍Cauchy导数定义外, 还介绍了与之等价的所谓Carathéodory 导数定义(定理 6.1.2). 这个定义完全避免了因自变量的增量 \({\Delta x}\) 与函数值的相应增量 \({\Delta f}\) 之比 \(\frac{\Delta f}{\Delta x}\) 而出现的 \(\frac{0}{0}\) 不定型极限,并刻画了函数图形的割线的斜率

以连续方式逼近函数在可导点处的斜率,以及函数在可导点处连续的内在性质. 导数的这个定义在证明函数的有关可导性定理时,例如证明复合函数的可导性、反函数的可导性等时显示了它的优越性。

另一点与传统分析教材不同的是, 在一元函数微积分部分, 我没有介绍函数的微分概念, 而是把它放在多元函数部分的映射微分中一起讲. 这样做, 教学上行得通, 因为只在积分及不定积分中用到了微分的记号,只要在那里作些技术上的处理和规定, 不定积分及定积分的计算可以避开函数微分这个概念.

鉴于函数的限定展开 (即通常的 Taylor 展开) 在函数的不定型极限中的计算、广义积分的收敛性判断及函数图形的研究中的重要应用, 我特别用一章来介绍函数的限定展开及其广义限定展开. 这部分内容相对传统分析教材来说, 内容份量比较重, 论述比较详尽。

关于函数的广义积分,为了避免概念与结论的重复叙述, 我把广义积分分为三类: \({\int }_{a}^{{b}^{ - }}f\left( x\right) {dx}\) 、 \({\int }_{{a}^{ + }}^{b}f\left( x\right) {dx}\) 、 \({\int }_{{a}^{ + }}^{{b}^{ - }}f\left( x\right) {dx}\) ,从而把有界区间上无界函数的广义积分与无界区间上函数的广义积分的定义与结论作了统一处理, 并把广义积分的 Abel 判别法与 Dirichlet 判别法的证明统一为一个一般性定理的证明, 从而使广义积分的叙述和定理证明更为简洁.

第二块的度量空间是为以后各章作准备的, 它为多元函数微积分奠定理论基础. 这一章的内容大多选自高年级的点集拓扑和泛函分析的有关概念和知识. 如何在低年级的分析教材中引进现代数学的概念和理论,对一、二年级大学生进行现代数学思维的训练, 并使得学生易于接受, 难度很大. 我从具体例子入手经过分析和抽象,引入度量空间、集合的开集、闭集、完备空间、 紧空间及连通空间等拓扑概念,然后又回到 \({\mathbf{R}}^{n}\) 空间上来. 这样就利用一般度量空间的抽象理论统一处理了经典微积分学中的二维、三维欧氏空间及二元、三元函数的极限与连续性, 使学生的认识得到了深化和提高。

第三块的级数部分,我特别强调了正项数项级数 的 \({\mathrm{d}}^{\prime }\mathrm{{Ale}}\) - mbert 判别法、 Cauchy 判别法、 Riemann 判别法、 积分判别法、Raabe判别法及Bertrand判别法的分析比较, 从而向学生解释了不存在能判断所有正项级数的散性的万能判别法的事实. 对于函数项序列与函数项级数, 根据它们的收敛性的共性, 我把它们的有关结论统一表述,以便学生进行对比学习. 此外,我还特别介绍了有重要应用价值的一致等度连续函数族的概念及实与复的Stone-Weierstrass稠密性定理.

关于用函数项级数来构造处处连续处处不可导函数的例子, 我没有采用通常的方法去构造 Weierstrass 函数, 而是采用了发表在美国数学月刊1991年第198卷第5期上由Katsuura构造的函数。这个函数的处处不可导性证明虽然比较长 (作为教材的例子也许不是很恰当),但它与Weierstrass函数比较, 不仅逼近函数的图形直观, 而且极限函数的图形也非常形象. 我选择这个函数的另一个目的是它的构造与当今研究十分活跃的新的数学分支 ——分形几何联系密切,通过引进这个函数,可以及时向学生介绍某些现代数学研究前沿的发展状况,丰富、扩大学生的知识面, 激发学生的学习积极性及创造性.

第四块的多元函数微积分部分, 与传统的分析教材比较, 从体系、结构、内容选材及写法上作了较大的改革尝试.

首先,在微分一章,比较详细地介绍了 \({\mathbf{R}}^{n}\) 空间中的映射的微分概念及其基本性质. 对重要的一类可微映射一一微分同胚, 除了严格证明了局部微分同胚定理 (隐函数定理)、逆射 定理 (反函数定理,外,还介绍了区域不变性定理及全局微分同胚定理.

为了在第21章介绍第二型曲线、曲面积分, 我用了一章来介绍微分形式的概念和性质. 微分形式可以采用交错重线性形式的方式来介绍,但这样需要介绍很多的有关重线性形式的知识,我采用纯规则的方法, 虽然几何意义差些, 但通过介绍外代数就比较快地引导出了微分形式的概念.

在重积分一章, 正如前面所指出的, 用阶梯函数序列对定义了 \({\mathbf{R}}^{n}\) 空间上的实值函数的Rimann可积性及其Riemann 积分, 并用积分定义了有界集的Jordan测度, 继而证明了集合的三个 Jordan可测特征性定理. 此外, 我还介绍了 Lebesgue 零测度集的概念, 并利用它给出了函数Rimann 可积的一个充分必要条件(即定理 19.3.7)。这一章,我用了很大篇幅来介绍重积分的各种计算方法, 即降维法 (Fubini定理、纤维法、截 面法)、对称法、变元替换法. 这些方法我都给予了严格证明. 特别是重积分的变元替换法, 我采用逐步化简的证明方法, 思路清晰, 论证严密, 学生易于分析和理解.

本书的最后两章用来介绍 \({\mathbf{R}}^{n}\) 空间的子流形及其定向、函数及微分形式沿子流形的积分,在充分研究了 \({\mathbf{R}}^{n}\) 空间的曲线、曲面的基础上,通过分析球面的参数表达式比较直观的引进了 \({\mathbf{R}}^{n}\) 空间的子流形的概念及子流形用坐标图及坐标图册表述的等价定义.

关于子流形的定向,我首先从分析 \({\mathbf{R}}^{n}\) 的 \(k\) 维子空间的定向 入手,利用子流形切空间的定向,引入了用坐标图册定义的以及用子流形切空间的连续定向族定义的子流形定向的两个等价定义. 此外还介绍了子流形的边缘子流形及其诱导定向的概念.

\(k\) 维子流形的Stokes公式是第21章的一个重要部分,利用它把经典的Newton-Leibniz公式、Green公式、Stokes 公式、 Ostrogradski 公式统一了起来. 作为Stokes公式的应用, 还证明了一个特殊的Brouwer 不动点定理及 \(n\) 维球体上 \({C}^{2}\) 类向量场的零点定理.

本书每节后面都选配了一定数量的练习题. 一类是基础性的,一类是拓广性的. 拓广性的题目都是通过介绍新概念、新结论, 培养学生独立分析问题和解决问题的能力的. 这类题目难度比较大,初次接触可以不必理会它,重点放在基础性题目的演练上.

本书共21章,如按每学期18周、每周讲授4学时计算(另每周安排 2-4 学时的习题课) 4 学期共288 学时可以讲完.

在教学过程中既要注意学生对理论问题的学习和钻研,也要注重对学生计算能力的培养, 还要帮助学生正确处理好学习本课程的基本理论与接受新的数学知识的关系,避免学生盲目进行脱离实际的超前学习.

还要指出一点的是,本人的愿望是想把书写成:提问清楚、 行文简洁、论证严密。但由于本人才薄识浅、水平有限、教学经验不足,很多地方可能事与愿违,甚至有谬误之处,诚请读者批评指正,不胜感激. 另外,本书的一些章节的安排只是我本人的考虑和想法, 不一定成熟和恰当, 教师在使用本教材时一定要根据自己的教学及学生的实际情况进行适当调整和取舍。

最后借本书出版之机向关心、支持我编写工作的武汉大学教务处、数学系及中法数学中心的领导、老师和同事们表示我的衷心感谢。吴厚心老师对我的工作始终给予了热情鼓励,并多次给予具体指导。为了使本书更趋完善,徐森林教授还多次给我来信商讨修改事宜, 他们为此书的出版倾注了一份难得的心力. 责任编辑文小西先生对书稿作了十分细致、认真的编辑加工工作,为提高本书的质量付出了大量的辛勤劳动,使本书增色不少. 在此一并表示我对他们的万分谢意.

在我编写本书的过程中,特别参阅了书末所列的国内外著作, 在此向著者们表示我的衷心感谢。



\begin{flushright}
	邹应\\
	1994 年 8 月于珞珈山
\end{flushright}

\newpage

\begin{center}
	\textbf{\heiti \large 提示}
\end{center}

致敬邹应教授.

在之前公众号发布的\textbf{数学分析新讲}的重排本及上次发出华东师范大学-数学分析重排本后,在\LaTeX 工作室公众号有很多人留言想要邹应老师的数学分析重排本，本次对上册进行重排,因邹老师数分内容很多，下册敬请期待.采用\LaTeX 工作室网站内复原的于品教授-数学分析的模板.本作品仅供排版练习使用，所有版权属于原作者.

{\heiti 注意本次重排对多行公式进行了校对于对齐，修改了章节等细节问题，虽然没有逐页核查是否有错误，但是基本上不影响阅读，几乎可以作为学习的书籍去使用，若有网友愿意继续改进非常感谢，未来更新待定ing....，接下来主要在于对例题定理等加环境，还有一些核对错误的问题，因为邹老师的书籍很早了流传的pdf也很差有些看不清.}



\newpage